\begin{ouverture}[statistique]{11}{Proportions et évolutions}{Partie 4 · Statistiques et probabilités}
{Pourcentages · Références · Coefficients multiplicateurs}
\objectif{Identifier le total de référence d’une proportion.}
\objectif{Calculer évolutions successives et réciproques.}
\objectif{Lire une information chiffrée de manière critique.}
\prerequis{Fractions et calcul décimal.}
\end{ouverture}
\pagestyle{otMain}\essentiel\fontsize{11}{15}\selectfont
\essSection{1}{Proportions et pourcentages}
\begin{definition}Si une partie contient $n$ éléments parmi $N>0$, sa proportion est $p=n/N$. L’écriture $25\,\%$ signifie $25/100=0{,}25$. Ainsi $n=pN$ et, si $p>0$, $N=n/p$.\end{definition}
\begin{propriete}Si $A$ est une partie de $B$, elle-même partie de $E$, la proportion de $A$ dans $E$ est le produit de la proportion de $A$ dans $B$ par celle de $B$ dans $E$.\end{propriete}
\begin{exemple}40\,\% d’une classe fait du sport ; parmi ces élèves, 30\,\% nagent. Les nageurs considérés représentent $0{,}40\times0{,}30=0{,}12$, soit 12\,\% de la classe.\end{exemple}
\begin{attention}Passer de 20\,\% à 25\,\% représente 5 points de pourcentage, mais une hausse relative de $(25-20)/20=25\,\%$. Toujours nommer le total ou la valeur de départ.\end{attention}
\begin{savoirfaire}[Comparer des groupes]On additionne des effectifs avant de calculer une proportion globale. Une moyenne simple de deux pourcentages n’est valable que si leurs populations ont le même effectif.\end{savoirfaire}
\clearpage
\essSection{2}{Évolutions multiplicatives}
\begin{definition}De $V_1>0$ à $V_2$, la variation absolue est $V_2-V_1$, le coefficient multiplicateur $C=V_2/V_1$, et le taux d’évolution $t=(V_2-V_1)/V_1=C-1$. Ainsi $V_2=V_1(1+t)$.\end{definition}
\begin{propriete}Pour deux évolutions successives, les coefficients se multiplient : $C=(1+t_1)(1+t_2)$. Le taux global vaut $C-1$. L’évolution réciproque a pour coefficient $1/C$ et pour taux $1/C-1$.\end{propriete}
\begin{exemple}Après une hausse de 20\,\% puis une baisse de 20\,\%, le coefficient vaut $1{,}2\times0{,}8=0{,}96$ : la baisse globale est de 4\,\%. Pour annuler une hausse de 20\,\%, la baisse nécessaire vaut $1-1/1{,}2=1/6$, soit environ 16,67\,\%.\end{exemple}
\begin{savoirfaire}[Tableur]Écrire un taux sous forme de pourcentage ou de décimal, puis une formule du type \code{=B2*(1+C2)} pour la valeur finale. Pour retrouver le départ, diviser par le coefficient ; ne pas soustraire le même pourcentage.\end{savoirfaire}
\otIntertitre{Interpréter}
Vérifier l’unité, les dates, le périmètre et la référence de comparaison. Une évolution d’une commune ne décrit pas automatiquement celle de son agglomération.
\clearpage\rubrique{Découverte}
\begin{blocLivre}{ex-11-1}
\begin{exercice}[1]{Calculer une proportion}
Parmi 240 élèves, 84 viennent à vélo. Donner leur proportion sous forme de fraction irréductible, décimale et de pourcentage.
\end{exercice}
\end{blocLivre}
\begin{blocLivre}{ex-11-2}
\begin{exercice}[1]{Appliquer un pourcentage}
Un prix de 80 euros bénéficie d’une remise de 15\,\%. Calculer la remise puis le prix payé.
\end{exercice}
\end{blocLivre}
\begin{blocLivre}{ex-11-3}
\begin{exercice}[1]{Retrouver le total}
54 adhérents représentent 18\,\% d’une association. Combien y a-t-il d’adhérents ?
\end{exercice}
\end{blocLivre}
\begin{blocLivre}{ex-11-4}
\begin{exercice}[1]{Du départ à l’arrivée}
Une quantité passe de 120 à 150. Calculer la variation absolue, le coefficient multiplicateur et le taux d’évolution.
\end{exercice}
\end{blocLivre}
\begin{blocLivre}{ex-11-5}
\begin{exercice}[1]{Une baisse}
Une consommation passe de 250 à 210 kWh. Calculer son taux d’évolution.
\end{exercice}
\end{blocLivre}

\clearpage\rubrique{Entraînement}
\begin{blocLivre}{ex-11-6}
\begin{exercice}[2]{Un pourcentage de pourcentage}
Dans un lycée de 800 élèves, 45\,\% suivent une option. Parmi eux, 40\,\% choisissent le théâtre. Combien d’élèves choisissent le théâtre, et quelle part du lycée représentent-ils ?
\label{fig-ex-11-6}
\figureGeo{proportions-emboitees}
\end{exercice}
\end{blocLivre}
\begin{blocLivre}{ex-11-7}
\begin{exercice}[2]{Deux évolutions successives}
Un prix de 200 euros augmente de 10\,\%, puis baisse de 15\,\%. Calculer le prix final et le taux global.
\end{exercice}
\end{blocLivre}
\begin{blocLivre}{ex-11-8}
\begin{exercice}[2]{L’évolution réciproque}
Une population augmente de 25\,\%. Quel pourcentage de baisse la ramènerait exactement à son effectif initial ?
\end{exercice}
\end{blocLivre}
\begin{blocLivre}{ex-11-9}
\begin{exercice}[2]{Points et pourcentage}
Un taux de réussite passe de 60\,\% à 72\,\%. Donner la hausse en points et la hausse relative.
\end{exercice}
\end{blocLivre}
\begin{blocLivre}{ex-11-10}
\begin{exercice}[2]{Une proportion globale}
Une classe de 20 élèves compte 80\,\% de réussites, une autre de 30 élèves en compte 60\,\%. Calculer le taux de réussite de l’ensemble.
\end{exercice}
\end{blocLivre}

\clearpage\rubrique{Problèmes}
\begin{blocLivre}{ex-11-11}
\begin{exercice}[3]{Retrouver un prix hors taxe}
Un prix TTC est de 144 euros avec une TVA de 20\,\%. Trouver le prix hors taxe et le montant de TVA.
\end{exercice}
\end{blocLivre}
\begin{blocLivre}{ex-11-12}
\begin{exercice}[3]{Une baisse suivie d’un retour}
Un appareil perd 30\,\% de sa valeur. De quel taux doit ensuite augmenter sa valeur pour retrouver celle du départ ? Donner un résultat exact puis un arrondi au dixième de pourcentage.
\end{exercice}
\end{blocLivre}
\begin{blocLivre}{ex-11-13}
\begin{exercice}[3]{Un seuil après plusieurs hausses}
Un capital de 500 euros est multiplié chaque année par 1,1. Après combien d’années dépasse-t-il 700 euros pour la première fois ? On autorise un tableau ou une boucle, sans logarithme.
\end{exercice}
\end{blocLivre}
\begin{blocLivre}{ex-11-14}
\begin{exercice}[3]{Des données de population}
L’Insee indique 327734 habitants pour la commune de Nantes et 267991 pour celle de Bordeaux au 1er janvier 2023. Calculer l’écart puis le pourcentage de plus de Nantes par rapport à Bordeaux, arrondi au dixième. \par\smallskip\legende{Source : Insee Focus 374, figure 3b, décembre 2025, \url{https://www.insee.fr/fr/statistiques/8682718}. Il s’agit des communes, pas des métropoles.}
\end{exercice}
\end{blocLivre}
\begin{blocLivre}{ex-11-15}
\begin{exercice}[3]{Un taux constant à retrouver}
Un prix augmente d’un même taux positif lors de deux hausses successives. Au total, il augmente de 21\,\%. Déterminer ce taux sans tâtonnement.
\end{exercice}
\end{blocLivre}

\clearpage
